\begin{titlepage}
    \centering
    \vspace*{2cm}
    
    % 主标题
    {\Huge \bfseries 电磁场作业集 \par}
    \vspace{0.8cm}
    
    % 副标题+课程信息
    {\Large \bfseries 考前复习 \par}
    \vspace{0.4cm}
    {\Large 课程代码：EE2003 \quad 课程名称：电磁场 \par}
    \vspace{0.6cm}
    
    \rule{\linewidth}{0.2mm}
    \vspace{1cm}
    
    % 麦克斯韦名言
    {\itshape \large 电和磁的实验中最明显的现象是，物体之间的相互作用。\par}
    \vspace{0.4cm}
    {\large —— 麦克斯韦 \par}
    \vspace{1.2cm}
    
    % 作者信息
    {\Large 作者：浮萍 \par}
    \vspace{0.3cm}
    {\Large 版本：1.0 \par}
    \vspace{0.3cm}
    {\Large 邮箱：3071974740@qq.com \par}
    \vspace{0.3cm}
    {\Large GitHub：github.com/Duckweed-yhb \par}
    \vspace{1cm}
    
    % 版权信息
    {\large \copyright\ 2026 浮萍 版权所有 \par}
    {\large 未经授权，不得复制、传播或修改 \par}
    \vspace{0.5cm}
    
    \vfill
    {\large 生成日期：\today \par}
\end{titlepage}

\section*{版本说明}
\noindent\textbf{版本：V1.0}
\begin{itemize}
    \item \textbf{配套课程：}课程代码\textbf{EE2003}，课程名称\textbf{电磁场}；
    \item \textbf{参考教材：}《工程电磁场导论》，冯慈璋、马西奎主编，西安交通大学出版社2000版，全书内容严格对标教材目录与期末考点；
    \item \textbf{覆盖范围：}全书五大章节：\textbf{第1章静电场、第2章恒定电场、第3章恒定磁场、第4章时变电磁场、第5章准静态电磁场}，对应教材$\S1\sim\S5$全部课内知识点；
    \item \textbf{内容编排：}
    \begin{itemize}
        \item 1.静电场：电场/电位、高斯定理、边值问题、有限差分法(实验内容)、镜像法、电容、静电能量与力（采用库仑定律+虚位移法求电场力）；
        \item 2.恒定电场：导电媒质电流、局外场强、基本方程、分界面衔接条件、恒定电场与静电场类比、电导；
        \item 3.恒定磁场：磁感应强度、安培环路定律、基本方程、磁矢位、标量磁位、磁场能量与力（洛伦兹公式+虚位移法求磁场力）；
        \item 4.时变电磁场：法拉第电磁感应+全电流定律、麦克斯韦方程组、动态位及其积分解、坡印廷矢量、正弦电磁场、电磁辐射；
        \item 5.准静态场：EQS电准静态/MQS磁准静态场、似稳条件、集肤效应、导体交流内阻抗；
    \end{itemize}
    \item 排版规范：关键公式统一封装\verb|formulaBox|环境，核心定义与考点名词全部\textbf{加粗}标注，各小节以\verb|\hrule|分隔，结构规整便于考前翻阅背诵；
    \item 内容取舍：剔除课外超纲拓展内容，只收录课堂核心考点、典型结论与工程实例，完全贴合课程期末考试大纲；
    \item \textbf{现存缺点：全篇仅文字与公式，无配套几何示意图、原理插图。}
\end{itemize}